%% kapa: exam=1 round=26 year=2015 period=1 form=A subject=civil_law subject_ko=민법 seq=4 number=4
%% source: Q-Net 감정평가사 기출문제 (cst00309) artlSeq=5138001 file=extracted/1008/5138001/제26회_감정평가사 1차_기출/제26회_감정평가사 1차_기출/1교시_감평_A형.pdf
%% conversion: kordoc 4.24.51 -> tools/gen1.py (auto); review=proofread-sonnet; flags=none
\begin{question}{4}
비법인사단 및 재단에 관한 설명으로 옳지 않은 것은? (다툼이 있으면 판례에 따름)
\begin{choices}
  \item 비법인사단의 대표자가 총유물의 관리ㆍ처분과 무관한 대외적 거래행위에 관하여 사원총회의 결의를 거치도록 한 정관 규정에 위반하여 그러한 거래행위를 한 경우, 상대방이 그와 같은 대표권 제한 사실을 알 수 없었다면 그 거래행위는 유효하다.
  \item 규약에 달리 정한 바가 없으면, 종중이 그 명의로 총유재산에 대한 보존행위로서 소송을 하기 위해서 종중총회의 결의를 거쳐야 하는 것은 아니다.
  \item 비법인사단에 대하여는 법인격을 전제로 하는 것을 제외하고는 사단법인에 관한 민법규정을 유추 적용한다.
  \item 매매계약에 의하여 부담하고 있는 채무의 존재를 인식하고 있다는 뜻을 표시함에 불과한 소멸시효 중단사유로서의 승인은 총유물의 관리ㆍ처분행위라고 볼 수 없다.
  \item 비법인재단의 경우에도 대표자가 있는 때에는 재단명의로 그 재단에 속하는 부동산의 등기를 할 수 있다.
\end{choices}
\end{question}
\answer{2}
